% Generated by tools/textbook/gen_credits.py from tools/assets/credits.json: do not edit.
\noindent{\footnotesize\textbf{Crediti delle immagini.} Le fotografie e gli schemi elencati qui provengono da Wikimedia Commons e conservano le proprie licenze; le altre figure sono state realizzate per questo libro. Nell'ordine dei nomi dei file: \fname{abacus\_chinese.jpg}: \emph{Abacus 5}, Loadmaster (David R. Tribble), CC BY-SA 3.0; \fname{abacus\_roman.jpg}: \emph{Number Stories, p51-illustration}, David Eugene Smith, Public domain; \fname{analytical\_engine.jpg}: \emph{Babbages Analytical Engine, 1834-1871. (9660574685)}, Science Museum London / Science and Society Picture Library, CC BY-SA 2.0; \fname{antikythera.jpg}: \emph{NAMA Machine d'Anticythère 1}, Marsyas, CC BY 2.5; \fname{antikythera\_plate.jpg}: \emph{Antikythera mechanism clockface, 1st-2nd century BC, Greece (model)}, Gts-tg, CC BY-SA 4.0; \fname{astrolabe.jpg}: \emph{Astrolabe planisférique closeup800x600x300}, Pom$^2$, CC BY-SA 3.0; \fname{c64.jpg}: \emph{Commodore-64-Computer-FL}, Evan-Amos, Public domain; \fname{cdrom\_drive.jpg}: \emph{CD Player focusing lens assembly}, Zim 256, CC BY-SA 3.0; \fname{chip\_74273.jpg}: \emph{Animation Technologies LifeView LR25 - Hitachi LS273-2523}, Raimond Spekking, CC BY-SA 4.0; \fname{colossus.jpg}: \emph{Dorothy Du Boisson}, unknown author (UK Government, 1943), Public domain; \fname{core\_memory.jpg}: \emph{HP 9101A extended memory unit.- magnetic core memory module - macro detail}, Mister rf, CC BY-SA 4.0; \fname{crt\_scope.jpg}: \emph{Tennis For Two on a DuMont Lab Oscilloscope Type 304-A}, Brookhaven National Laboratory, Public domain; \fname{difference\_engine.jpg}: \emph{Babbage Difference Engine}, geni, CC BY-SA 4.0; \fname{dvd\_drive.jpg}: \emph{Dvd-2010a hg}, Hannes Grobe, CC BY-SA 4.0; \fname{eniac.png}: \emph{Glen Beck and Betty Snyder program the ENIAC in building 328 at the Ballistic Research Laboratory}, U.S. Army, Public domain; \fname{enigma.jpg}: \emph{Enigma I cipher machine Museum of Communication Berlin cropped}, Museum für Kommunikation Berlin, CC BY 4.0; \fname{eprom.jpg}: \emph{Pair of BIOS chips}, Jud McCranie, CC BY-SA 4.0; \fname{hard\_disk.jpg}: \emph{Open hard-drive}, Zzubnik, Public domain; \fname{laser\_printer.jpg}: \emph{Laser printer-Writing}, DMahalko, traced by Stannered, CC BY 3.0; \fname{lcd\_layers.jpg}: \emph{TN-LCD-schematic-MS-208kB}, see the file page, CC BY-SA 3.0; \fname{lcd\_pixels.jpg}: \emph{Subpixel rendering LCD photo 3e composite}, ALexL33, CC BY-SA 3.0; \fname{leibniz\_reckoner.jpg}: \emph{Leibniz-Rechenmaschine hg}, Hannes Grobe, CC BY 3.0; \fname{mac\_classic.jpg}: \emph{Macintosh 128k transparency}, Grm wnr, CC BY-SA 3.0; \fname{mainframe.jpg}: \emph{Dept of Agriculture Computer Room}, Warren K. Leffler (U.S. News \& World Report, Library of Congress), Public domain; \fname{minicomputer.jpg}: \emph{DEC PDP-11 Minicomputer in extended variant (edited)}, Loz Pycock, CC BY-SA 2.0; \fname{pascaline.jpg}: \emph{Arts et Metiers Pascaline dsc03869}, David.Monniaux, CC BY-SA 3.0; \fname{pcb.jpg}: \emph{Hand Etched PCB}, Tpdwkouaa, CC BY-SA 4.0; \fname{pit\_land.jpg}: \emph{CD-ROM pressed pits}, Tycho, CC0; \fname{punched\_cards.jpg}: \emph{A Jacquard loom showing information punchcards, National Museum of Scotland}, Stephencdickson, CC BY-SA 4.0; \fname{ram\_modules.jpg}: \emph{Platinen-001 hg}, Hannes Grobe, CC BY-SA 4.0; \fname{ram\_on\_board.jpg}: \emph{IBM PS2 MCA Model 55 SX, power supply, memory, riser base}, Dale Mahalko, CC BY-SA 3.0; \fname{ssd.jpg}: \emph{Vertex 2 Solid State Drive by OCZ-top oblique PNr\textdegree{}0307}, D-Kuru, CC BY-SA 3.0 AT; \fname{storage\_racks.jpg}: \emph{StorageTek tape library (51993015872)}, Jorge Franganillo, CC BY 2.0; \fname{workstation\_board.jpg}: \emph{IBM PC Server 320 motherboard front}, deksor, CC BY-SA 4.0; \fname{zuse\_z3.jpg}: \emph{Z3 Deutsches Museum b}, Venusianer, CC BY-SA 3.0.\par}
